\section{A parity-preserving theta approximation that cannot improve the zero set}
\label{sec:gaussianrepair}
A natural response to the endpoint defect of
Theorem~\ref{thm:falsecutoff} is to smooth each finite
theta-mode approximation. We show exactly what such
smoothing achieves and why it does not constitute
the missing zero-location argument.

\begin{theorem}[Positive analytic parity repair and zero invariance]
\label{thm:parityrepair}
Let $M\ge1$ and
$f_M(r)=\sum_{m=1}^{M}\phi_m(r)>0$ for $r\ge0$,
as in Section~\ref{sec:squarelattice}. Put
$P_M(u)=f_M(|u|)$ and define, for $\sigma>0$,
\begin{equation}
K_{M,\sigma}(u)=
\frac{1}{\sqrt{2\pi}\sigma}
\int_{\mathbb R}
e^{-(u-v)^2/(2\sigma^2)}P_M(v)\,dv .
\label{eq:parityrepairK}
\end{equation}
Then $K_{M,\sigma}$ is positive, even, real
analytic, rapidly decreasing, and has all odd
derivatives zero at $u=0$.
Its Fourier transform is the entire function
\begin{equation}
\boxed{
F_{M,\sigma}(z)=e^{-\sigma^2z^2/2}F_M(z).
}
\label{eq:parityrepairF}
\end{equation}
In particular $F_{M,\sigma}$ and $F_M$
have \emph{exactly the same complex zeros
with the same multiplicities}.
Their first generalized Laguerre expressions
satisfy the pointwise exact identity
\begin{equation}
\boxed{
L_1[F_{M,\sigma}](x)
=e^{-\sigma^2x^2}
\bigl(L_1[F_M](x)+\sigma^2F_M(x)^2\bigr).
}
\label{eq:parityrepairL}
\end{equation}
For every fixed $n\ge0$ and $\sigma>0$,
as real $|x|\to\infty$,
\begin{equation}
\boxed{
L_n[F_{M,\sigma}](x)
=e^{-\sigma^2x^2}\frac{4D_M^2}{x^4}
\left(\frac{\sigma^{2n}}{n!}
+O_{M,\sigma,n}(x^{-2})\right),
}
\label{eq:parityrepairasym}
\end{equation}
where $D_M=f_M'(0)>0$.
Thus \emph{every fixed Laguerre order} of
the smoothed cutoff is eventually strictly
positive on the real axis, despite its
unchanged complex zero set.
\end{theorem}

\begin{proof}
The integrand in \eqref{eq:parityrepairK} is
strictly positive, and convolution of an
even integrable function with an even
Gaussian is even. Gaussian differentiation
under the integral proves real analyticity
and Schwartz decay; all odd derivatives
vanish at zero by parity. The Fourier
convolution theorem initially on the real
axis, followed by holomorphic continuation,
gives \eqref{eq:parityrepairF}. Since
$e^{-\sigma^2z^2/2}$ never vanishes,
its zeros and their multiplicities are
those of $F_M$.

For $g(x)=e^{-\sigma^2x^2/2}$,
\[
\frac{g'(x)}{g(x)}=-\sigma^2x,\qquad
\frac{g''(x)}{g(x)}=\sigma^4x^2-\sigma^2.
\]
Expanding $(g'F+gF')^2-(gF)(g''F+
2g'F'+gF'')$ and cancelling cross terms
proves \eqref{eq:parityrepairL}.

The repeated integration-by-parts
expansion in Theorem~\ref{thm:falsecutoff}
extends uniformly to a fixed complex strip
$|\operatorname{Im}z|\le R$:
\[
F_M(z)=-2D_Mz^{-2}+O_{M,R}(z^{-4}),
\quad |\operatorname{Re}z|\to\infty.
\]
Uniform strip bounds also justify
termwise extraction of finitely many
coefficients in the local complex
variable $y$ by Cauchy's formula.
Indeed
\begin{align*}
F_{M,\sigma}(x+iy)F_{M,\sigma}(x-iy)
={}&e^{-\sigma^2x^2}e^{\sigma^2y^2}\\
&\cdot F_M(x+iy)F_M(x-iy).
\end{align*}
For bounded complex $y$ the last product
is $4D_M^2x^{-4}+O_{M,R}(x^{-6})$.
The coefficient of $y^{2n}$ in
$e^{\sigma^2y^2}$ is $\sigma^{2n}/n!$.
This gives \eqref{eq:parityrepairasym}.
Its leading coefficient is strictly
positive for every fixed $n$ and
$\sigma>0$, completing the proof.
\end{proof}

\begin{theorem}[Quantitative uniform approximation
with unchanged cutoff zeros]
\label{thm:parityuniform}
Let
\[
E_M=3q_{M+1}e^{-q_{M+1}},\qquad
J_M=\int_0^\infty|f_M''(r)|\,dr<\infty.
\]
Then for all real $x$,
\begin{equation}
\boxed{
|F(x)-F_{M,\sigma}(x)|
\le E_M+\sigma^2(D_M+J_M).
}
\label{eq:uniformgauss}
\end{equation}
In particular, choosing
$\sigma_M^2=E_M/(1+D_M+J_M)$ gives
\[
\sup_{x\in\mathbb R}
|F(x)-F_{M,\sigma_M}(x)|<2E_M
\longrightarrow0.
\]
Thus positive, even, real-analytic
approximants with exact odd-jet parity
converge \emph{uniformly on the full real
Fourier axis} to the genuine Riemann
$\Xi$ transform, while each retains
the exact complex zeros of its
unsmoothed finite theta-mode cutoff.
\end{theorem}
\begin{proof}
Theorem~\ref{thm:uniformFourier} with
$k=0$ gives
$\sup_x|F-F_M|\le E_M$.
Twice integrating by parts, for $x\ne0$,
\[
x^2F_M(x)
=-2D_M-2\int_0^\infty
f_M''(r)\cos(xr)\,dr,
\]
so $x^2|F_M(x)|\le2(D_M+J_M)$.
The elementary inequality
$0\le1-e^{-t}\le t$ for $t\ge0$ gives
\[
|F_M(x)-F_{M,\sigma}(x)|
\le\frac{\sigma^2x^2}{2}|F_M(x)|
\le\sigma^2(D_M+J_M).
\]
Continuity covers $x=0$, proving
\eqref{eq:uniformgauss}.
The displayed choice of $\sigma_M$
makes the second term strictly
less than $E_M$, while $E_M\to0$
superexponentially in $M^2$.
\end{proof}

\paragraph{Independent tests.}
For $M=1$, $\sigma=0.35$ and real
$x\in\{39.5,60,100\}$, direct
incomplete-gamma differentiation
of $F_M$ verifies
$L_1[F_M](x)<0$ but
\[
L_1[F_M](x)+\sigma^2F_M(x)^2>0
\]
at all three points. This is a
numerical cross-check, whereas
\eqref{eq:parityrepairL} and the
eventual-positive theorem are exact
analytic identities.
\texttt{certificates/gaussian\_parity\_repair\_exact.py}
checks the polynomial/inequality
gates in rational arithmetic.
The Lean module
\texttt{XiKernel/GaussianParityRepair.lean}
formalizes the first Laguerre
product law and preservation of
complex zeros under a Gaussian
multiplier. It does not formalize
all convolution integrals,
Fourier asymptotics or the
$\Xi$-specific uniform bound
end-to-end.

\paragraph{Proof boundary and prior art.}
Gaussian smoothing is a classical
operation, and finite Fourier
transforms of Jacobi-theta kernels
with nonreal zeros were studied
by Csordas \cite{CsTrunc}.
The point here is a falsification
of a tempting \emph{new proof route}:
even global uniform real-axis
approximation, positivity,
analytic parity repair, and
eventual positivity at every
fixed Laguerre order cannot
certify the real-zero property
without additional genuinely
modular, global-in-order
Fourier information.
This theorem does not assert
that the true $\Xi$ transform
has any nonreal zero, and
does not resolve RH.
